\begin{thebibliography}{74}

\bibitem[1]{ref1}
Eric Allender and Bireswar Das. Zero knowledge and circuit minimization. Inf. Comput., 256:2–8, 2017. doi:10.1016/j.ic.2017.04.004.

\bibitem[2]{ref2}
Vikraman Arvind and Jacobo Torán. Isomorphism testing: Perspective and open problems. Bulletin of the EATCS, 86:66–84, 2005.

\bibitem[3]{ref3}
László Babai. Graph isomorphism in quasipolynomial time [extended abstract]. In Proceedings of the 48th Annual ACM SIGACT Symposium on Theory of Computing, STOC 2016, pages 684–697, 2016. arXiv:1512.03547 [cs.DS] version 2. doi:10.1145/2897518.2897542.

\bibitem[4]{ref4}
László Babai, Paolo Codenotti, Joshua A. Grochow, and Youming Qiao. Code equivalence and group isomorphism. In Proceedings of the Twenty-Second Annual ACM–SIAM Symposium on Discrete Algorithms (SODA11), pages 1395–1408, Philadelphia, PA, 2011. SIAM. doi: 10.1137/1.9781611973082.107.

\bibitem[5]{ref5}
László Babai, Paolo Codenotti, and Youming Qiao. Polynomial-time isomorphism test for groups with no abelian normal subgroups - (extended abstract). In Automata, Languages, and Programming - 39th International Colloquium, ICALP 2012, Proceedings, Part I, pages 51–62, 2012. doi:10.1007/978-3-642-31594-7\_5.

\bibitem[6]{ref6}
László Babai and Youming Qiao. Polynomial-time isomorphism test for groups with Abelian Sylow towers. In 29th STACS, pages 453–464. Springer LNCS 6651, 2012. doi:10.4230/ LIPIcs.STACS.2012.453.

\bibitem[7]{ref7}
Reinhold Baer. Groups with abelian central quotient group. Trans. AMS, 44(3):357–386, 1938. doi:10.1090/S0002-9947-1938-1501972-1.

\bibitem[8]{ref8}
Mihir Bellare and Shafi Goldwasser. The complexity of decision versus search. SIAM J. Comput., 23(1):97–119, 1994. doi:10.1137/S0097539792228289.

\bibitem[9]{ref9}
Hans Ulrich Besche and Bettina Eick. Construction of finite groups. J. Symb. Comput., 27(4):387–404, 1999. doi:10.1006/jsco.1998.0258.

\bibitem[10]{ref10}
Hans Ulrich Besche, Bettina Eick, and E.A. O’Brien. A millennium project: Constructing small groups. Intern. J. Alg. and Comput, 12:623–644, 2002. doi:10.1142/S0218196702001115.

\bibitem[11]{ref11}
Anton Betten, Michael Braun, Harald Fripertinger, Adalbert Kerber, Axel Kohnert, and Alfred Wassermann. Error-correcting linear codes: Classification by isometry and applications, volume 18. Springer Science and Business Media, 2006.

\bibitem[12]{ref12}
Peter A. Brooksbank, Joshua A. Grochow, Yinan Li, Youming Qiao, and James B. Wilson. Incorporating Weisfeiler–Leman into algorithms for group isomorphism. arXiv:1905.02518 [cs.CC], 2019.

\bibitem[13]{ref13}
Peter A. Brooksbank, Yinan Li, Youming Qiao, and James B. Wilson. Improved algorithms for alternating matrix space isometry: From theory to practice. In Fabrizio Grandoni, Grzegorz Herman, and Peter Sanders, editors, 28th Annual European Symposium on Algorithms, ESA 2020, September 7-9, 2020, Pisa, Italy (Virtual Conference), volume 173 of LIPIcs, pages 26:1–26:15. Schloss Dagstuhl - Leibniz-Zentrum für Informatik, 2020. doi:10.4230/LIPIcs. ESA.2020.26.

\bibitem[14]{ref14}
Peter A. Brooksbank and Eugene M. Luks. Testing isomorphism of modules. J. Algebra, 320(11):4020–4029, 2008. doi:10.1016/j.jalgebra.2008.07.014.

\bibitem[15]{ref15}
Peter A. Brooksbank, Joshua Maglione, and James B. Wilson. A fast isomorphism test for groups whose Lie algebra has genus 2. J. Algebra, 473:545–590, 2017. doi:10.1016/j. jalgebra.2016.12.007.

\bibitem[16]{ref16}
Peter A. Brooksbank and E. A. O’Brien. Constructing the group preserving a system of forms. Internat. J. Algebra Comput., 18(2):227–241, 2008. doi:10.1142/S021819670800441X.

\bibitem[17]{ref17}
John J. Cannon and Derek F. Holt. Automorphism group computation and isomorphism testing in finite groups. J. Symbolic Comput., 35(3):241–267, 2003. doi:10.1016/S0747-7171(02) 00133-5.

\bibitem[18]{ref18}
Xi Chen, Xiaotie Deng, and Shang-Hua Teng. Settling the complexity of computing two-player Nash equilibria. J. ACM, 56(3):Art. 14, 57, 2009. doi:10.1145/1516512.1516516.

\bibitem[19]{ref19}
Alexander Chistov, Gábor Ivanyos, and Marek Karpinski. Polynomial time algorithms for modules over finite dimensional algebras. In Proceedings of the 1997 International Symposium on Symbolic and Algebraic Computation, ISSAC ’97, pages 68–74. ACM, 1997. doi:10.1145/ 258726.258751.

\bibitem[20]{ref20}
Serena Cicalò, Willem A. de Graaf, and Michael Vaughan-Lee. An effective version of the Lazard correspondence. J. Algebra, 352(1):430–450, 2012. doi:10.1016/j.jalgebra.2011.11.031.

\bibitem[21]{ref21}
W.A. de Graaf. Lie Algebras: Theory and Algorithms, volume 56 of North-Holland Mathemat- ical Library. Elsevier Science, 2000.

\bibitem[22]{ref22}
Holger Dell, Thore Husfeldt, Dániel Marx, Nina Taslaman, and Martin Wahlén. Exponential time complexity of the permanent and the Tutte polynomial. ACM Trans. Algorithms, 10(4):Art. 21, 32, 2014. doi:10.1145/2635812.

\bibitem[23]{ref23}
Sean Eberhard and Stefan-C. Virchow. Random generation of the special linear group. Transactions of the American Mathematical Society, page 1, 2020. doi:10.1090/tran/8009.

\bibitem[24]{ref24}
Wayne Eberly and Mark Giesbrecht. Efficient decomposition of associative algebras over finite fields. Journal of Symbolic Computation, 29(3):441–458, 2000. doi:10.1006/jsco.1999.0308.

\bibitem[25]{ref25}
Bettina Eick, C. R. Leedham-Green, and E. A. O’Brien. Constructing automorphism groups of p-groups. Comm. Algebra, 30(5):2271–2295, 2002. doi:10.1081/AGB-120003468.

\bibitem[26]{ref26}
V. Felsch and J. Neubüser. On a programme for the determination of the automorphism group of a finite group. In Pergamon J. Leech, editor, Computational Problems in Abstract Algebra (Proceedings of a Conference on Computational Problems in Algebra, Oxford, 1967), pages 59–60, Oxford, 1970.

\bibitem[27]{ref27}
Katalin Friedl and Lajos Rónyai. Polynomial time solutions of some problems in computational algebra. In Robert Sedgewick, editor, Proceedings of the 17th Annual ACM Symposium on Theory of Computing, May 6-8, 1985, Providence, Rhode Island, USA, pages 153–162. ACM, 1985. doi:10.1145/22145.22162.

\bibitem[28]{ref28}
Vyacheslav Futorny, Joshua A. Grochow, and Vladimir V. Sergeichuk. Wildness for tensors. Lin. Alg. Appl., 566:212–244, 2019. doi:10.1016/j.laa.2018.12.022.

\bibitem[29]{ref29}
Joshua A. Grochow. Answer to “what is the hardest instance for the group isomorph- ism problem?”. Theoretical Computer Science Stack Exchange. URL: https://cstheory. stackexchange.com/a/42551/129.

\bibitem[30]{ref30}
Joshua A. Grochow and Youming Qiao. Polynomial-time isomorphism test of groups that are tame extensions - (extended abstract). In Algorithms and Computation - 26th International Symposium, ISAAC 2015, Nagoya, Japan, December 9-11, 2015, Proceedings, pages 578–589, 2015. doi:10.1007/978-3-662-48971-0\_49.

\bibitem[31]{ref31}
Joshua A. Grochow and Youming Qiao. Algorithms for group isomorphism via group extensions and cohomology. SIAM J. Comput., 46(4):1153–1216, 2017. Preliminary version in IEEE Conference on Computational Complexity (CCC) 2014 (DOI:10.1109/CCC.2014.19). Also available as arXiv:1309.1776 [cs.DS] and ECCC Technical Report TR13-123. doi:10.1137/ 15M1009767.

\bibitem[32]{ref32}
Joshua A. Grochow and Youming Qiao. On the complexity of isomorphism problems for tensors, groups, and polynomials I: Tensor Isomorphism-completeness. In ITCS, page to appear, 2021. arXiv:1907.00309.

\bibitem[33]{ref33}
Martin Grohe and Pascal Schweitzer. The graph isomorphism problem. Commun. ACM, 63(11):128–134, 2020. doi:10.1145/3372123.

\bibitem[34]{ref34}
Xiaoyu He and Youming Qiao. On the Baer–Lovász–Tutte construction of groups from graphs: isomorphism types and homomorphism notions. arXiv:2003.07200 [math.CO], 2020.

\bibitem[35]{ref35}
Russell Impagliazzo and Ramamohan Paturi. On the complexity of k-SAT. J. Comput. System Sci., 62(2):367–375, 2001. Special issue on the Fourteenth Annual IEEE Conference on Computational Complexity (Atlanta, GA, 1999). doi:10.1006/jcss.2000.1727.

\bibitem[36]{ref36}
Gábor Ivanyos. Fast randomized algorithms for the structure of matrix algebras over finite fields. In Proceedings of the 2000 international symposium on Symbolic and algebraic computation, pages 175–183. ACM, 2000. doi:10.1145/345542.345620.

\bibitem[37]{ref37}
Gábor Ivanyos, Marek Karpinski, and Nitin Saxena. Deterministic polynomial time algorithms for matrix completion problems. SIAM J. Comput., 39(8):3736–3751, 2010. doi:10.1137/ 090781231.

\bibitem[38]{ref38}
Gábor Ivanyos and Lajos Rónyai. Computations in associative and Lie algebras. In Some tapas of computer algebra, pages 91–120. Springer, 1999. doi:10.1007/978-3-662-03891-8\_5.

\bibitem[39]{ref39}
Zhengfeng Ji, Youming Qiao, Fang Song, and Aaram Yun. General linear group action on tensors: A candidate for post-quantum cryptography. In Dennis Hofheinz and Alon Rosen, editors, Theory of Cryptography - 17th International Conference, TCC 2019, Nuremberg, Germany, December 1-5, 2019, Proceedings, Part I, volume 11891 of Lecture Notes in Computer Science, pages 251–281. Springer, 2019. Preprint arXiv:1906.04330 [cs.CR]. doi:10.1007/ 978-3-030-36030-6\_11.

\bibitem[40]{ref40}
William M. Kantor. Some topics in asymptotic group theory. Groups, Combinatorics and Geometry (Durham, pages 403–421, 1990.

\bibitem[41]{ref41}
William M Kantor and Alexander Lubotzky. The probability of generating a finite classical group. Geometriae Dedicata, 36(1):67–87, 1990.

\bibitem[42]{ref42}
Neeraj Kayal and Timur Nezhmetdinov. Factoring groups efficiently. In Susanne Albers, Alberto Marchetti-Spaccamela, Yossi Matias, Sotiris E. Nikoletseas, and Wolfgang Thomas, editors, Automata, Languages and Programming, 36th International Colloquium, ICALP 2009, Rhodes, Greece, July 5-12, 2009, Proceedings, Part I, volume 5555 of Lecture Notes in Computer Science, pages 585–596. Springer, 2009. Preprint ECCC Tech. Report TR08-074. doi:10.1007/978-3-642-02927-1\_49.

\bibitem[43]{ref43}
E. I. Khukhro. p-automorphisms of finite p-groups, volume 246 of London Mathematical Society Lecture Note Series. Cambridge University Press, Cambridge, 1998. doi:10.1017/ CBO9780511526008.

\bibitem[44]{ref44}
Johannes Köbler, Uwe Schöning, and Jacobo Torán. The graph isomorphism problem: its structural complexity. Birkhauser Verlag, Basel, Switzerland, Switzerland, 1993. doi:10.1007/ 978-1-4612-0333-9.

\bibitem[45]{ref45}
Tamara G Kolda and Brett W Bader. Tensor decompositions and applications. SIAM review, 51(3):455–500, 2009. doi:10.1137/07070111X.

\bibitem[46]{ref46}
Michel Lazard. Sur les groupes nilpotents et les anneaux de Lie. Ann. Sci. Ecole Norm. Sup. (3), 71:101–190, 1954. doi:10.24033/asens.1021.

\bibitem[47]{ref47}
François Le Gall. Efficient isomorphism testing for a class of group extensions. In Proc. 26th STACS, pages 625–636, 2009. doi:10.4230/LIPIcs.STACS.2009.1830.

\bibitem[48]{ref48}
Mark L. Lewis and James B. Wilson. Isomorphism in expanding families of indistinguishable groups. Groups Complex. Cryptol., 4(1):73–110, 2012. doi:10.1515/gcc-2012-0008.

\bibitem[49]{ref49}
Yinan Li and Youming Qiao. Linear algebraic analogues of the graph isomorphism problem and the Erdős–Rényi model. In Chris Umans, editor, 58th IEEE Annual Symposium on Foundations of Computer Science, FOCS 2017, pages 463–474. IEEE Computer Society, 2017. doi:10.1109/FOCS.2017.49.

\bibitem[50]{ref50}
Richard J. Lipton, Lawrence Snyder, and Yechezkel Zalcstein. The complexity of word and isomorphism problems for finite groups. Yale University Department of Computer Science Research Report \# 91, 1977. URL: https://apps.dtic.mil/dtic/tr/fulltext/u2/a053246. pdf.

\bibitem[51]{ref51}
Eugene M. Luks. Computing in solvable matrix groups. In FOCS 1992, 33rd Annual Symposium on Foundations of Computer Science, pages 111–120. IEEE Computer Society, 1992. doi:10.1109/SFCS.1992.267813.

\bibitem[52]{ref52}
Eugene M. Luks. Permutation groups and polynomial-time computation. In Groups and computation (New Brunswick, NJ, 1991), volume 11 of DIMACS Ser. Discrete Math. Theoret. Comput. Sci., pages 139–175. Amer. Math. Soc., Providence, RI, 1993.

\bibitem[53]{ref53}
Eugene M. Luks. Hypergraph isomorphism and structural equivalence of boolean functions. In Proceedings of the Thirty-First Annual ACM Symposium on Theory of Computing, May 1-4, 1999, Atlanta, Georgia, USA, pages 652–658, 1999. doi:10.1145/301250.301427.

\bibitem[54]{ref54}
Rudolf Mathon. A note on the graph isomorphism counting problem. Information Processing Letters, 8(3):131–136, 1979.

\bibitem[55]{ref55}
Brendan D. McKay. Practical graph isomorphism. Congr. Numer., pages 45–87, 1980.

\bibitem[56]{ref56}
Brendan D. McKay and Adolfo Piperno. Practical graph isomorphism, II. Journal of Symbolic Computation, 60(0):94–112, 2014. doi:10.1016/j.jsc.2013.09.003.

\bibitem[57]{ref57}
Alan H. Mekler. Stability of nilpotent groups of class 2 and prime exponent. The Journal of Symbolic Logic, 46(4):781–788, 1981.

\bibitem[58]{ref58}
Gary L. Miller. On the $n^{\log n}$ isomorphism technique (a preliminary report). In STOC, pages 51–58. ACM, 1978. doi:10.1145/800133.804331.

\bibitem[59]{ref59}
Takunari Miyazaki. Luks’s reduction of graph isomorphism to code equivalence. Comment to E. W. Clark, https://groups.google.com/forum/\#!msg/sci.math.research/puZxGj9HXKI/ CeyH2yyyNFUJ, 1996.

\bibitem[60]{ref60}
Vipul Naik. Lazard correspondence up to isoclinism. PhD thesis, The University of Chicago, 2013. URL: https://vipulnaik.com/thesis/.

\bibitem[61]{ref61}
Jacques Patarin. Hidden fields equations (HFE) and isomorphisms of polynomials (IP): two new families of asymmetric algorithms. In Advances in Cryptology - EUROCRYPT ’96, Inter- national Conference on the Theory and Application of Cryptographic Techniques, Saragossa, Spain, May 12-16, 1996, Proceeding, pages 33–48, 1996. doi:10.1007/3-540-68339-9\_4.

\bibitem[62]{ref62}
Erez Petrank and Ron M. Roth. Is code equivalence easy to decide? IEEE Trans. Inf. Theory, 43(5):1602–1604, 1997. doi:10.1109/18.623157.

\bibitem[63]{ref63}
Youming Qiao, Jayalal M. N. Sarma, and Bangsheng Tang. On isomorphism testing of groups with normal Hall subgroups. In Proc. 28th STACS, pages 567–578, 2011. doi: 10.4230/LIPIcs.STACS.2011.567.

\bibitem[64]{ref64}
Lajos Rónyai. Computing the structure of finite algebras. J. Symb. Comput., 9(3):355–373, 1990. doi:10.1016/S0747-7171(08)80017-X.

\bibitem[65]{ref65}
David J. Rosenbaum. Bidirectional collision detection and faster deterministic isomorphism testing. arXiv preprint arXiv:1304.3935 [cs.DS], 2013.

\bibitem[66]{ref66}
David J. Rosenbaum. Breaking the $n^{\log n}$ barrier for solvable-group isomorphism. In Proceedings of the Twenty-Fourth Annual ACM-SIAM Symposium on Discrete Algorithms, pages 1054–1073. SIAM, 2013. Preprint arXiv:1205.0642 [cs.DS]. doi:10.1137/1.9781611973105.76.

\bibitem[67]{ref67}
Nicolas Sendrier and Dimitris E. Simos. The hardness of code equivalence over $\mathbb F_q$ and its application to code-based cryptography. In International Workshop on Post-Quantum Cryptography, pages 203–216. Springer, 2013.

\bibitem[68]{ref68}
Seinosuke Toda. PP is as hard as the polynomial-time hierarchy. SIAM J. Comput., 20(5):865– 877, 1991. doi:10.1137/0220053.

\bibitem[69]{ref69}
Leslie G. Valiant. Relative complexity of checking and evaluating. Information Processing Lett., 5(1):20–23, 1976/77. doi:10.1016/0020-0190(76)90097-1.

\bibitem[70]{ref70}
James Wilson. 2014 conference on Groups, Computation, and Geometry at Colorado State University, co-organized by P. Brooksbank, A. Hulpke, T. Penttila, J. Wilson, and W. Kantor. Personal communication, 2014.

\bibitem[71]{ref71}
James B. Wilson. Decomposing p-groups via Jordan algebras. J. Algebra, 322(8):2642–2679, 2009. doi:10.1016/j.jalgebra.2009.07.029.

\bibitem[72]{ref72}
James B. Wilson. Finding direct product decompositions in polynomial time. arXiv:1005.0548 [math.GR], 2010.

\bibitem[73]{ref73}
James B. Wilson. Existence, algorithms, and asymptotics of direct product decompositions, I. Groups Complex. Cryptol., 4(1):33–72, 2012. doi:10.1515/gcc-2012-0007.

\bibitem[74]{ref74}
V. N. Zemlyachenko, N. M. Korneenko, and R. I. Tyshkevich. Graph isomorphism problem. J. Soviet Math., 29(4):1426–1481, May 1985. doi:10.1007/BF02104746.
\end{thebibliography}
